% !TEX root = ../DoAn.tex
\documentclass[../DoAn.tex]{subfiles}
\begin{document}

\section{Đặt vấn đề}
\label{section:1.1}

% Trong thập kỷ qua, thương mại điện tử (TMĐT) tại Việt Nam đã chứng kiến tốc độ tăng trưởng
% vượt bậc. Theo báo cáo của Hiệp hội Thương mại điện tử Việt Nam (VECOM) năm 2024, doanh thu
% TMĐT B2C ước đạt 25 tỷ USD, tăng trưởng bình quân trên 25\%/năm trong giai đoạn 2020--2024.
% Các sàn giao dịch lớn như Shopee, Lazada, Tiki đang phục vụ hàng chục triệu người dùng đồng
% thời, đặc biệt trong các sự kiện Flash Sale và Black Friday --- nơi lưu lượng giao dịch có
% thể tăng đột biến gấp 10--50 lần so với ngày thường.

% Thực tế vận hành cho thấy các hệ thống TMĐT xây dựng theo kiến trúc Monolith truyền thống
% gặp nhiều hạn chế nghiêm trọng khi đối mặt với bài toán quy mô lớn. Thứ nhất, toàn bộ hệ
% thống phải được mở rộng đồng thời dù chỉ một module đang chịu tải cao, dẫn đến lãng phí tài
% nguyên và chi phí vận hành. Thứ hai, mọi thay đổi dù nhỏ đều yêu cầu triển khai lại toàn bộ
% ứng dụng, làm chậm tốc độ phát hành tính năng mới và tăng rủi ro khi cập nhật. Thứ ba, một
% lỗi cục bộ có thể kéo sập toàn bộ hệ thống do thiếu cơ chế cách ly sự cố (fault isolation).
% Đối với các nghiệp vụ đặc thù như flash sale --- nơi hàng chục nghìn người dùng đồng thời
% cạnh tranh một số lượng hàng giới hạn trong vài phút --- Monolith không thể đáp ứng được yêu
% cầu về khả năng mở rộng chọn lọc và xử lý đồng thời cao.

% Kiến trúc microservices ra đời như một giải pháp cho những hạn chế nêu trên. Tuy nhiên, việc
% thiết kế và triển khai một hệ thống microservices hoàn chỉnh đòi hỏi nắm vững nhiều mẫu thiết
% kế phức tạp: từ Saga Pattern cho giao dịch phân tán, Transactional Outbox cho truyền thông
% điệp tin cậy, đến Idempotent Consumer cho xử lý phân phối at-least-once. Đây là những vấn đề
% được đề cập rộng rãi trong tài liệu học thuật nhưng ít dự án thực tế tại Việt Nam triển khai
% một cách đầy đủ và có hệ thống. Bên cạnh đó, việc tích hợp các thành phần hạ tầng như
% service discovery, API gateway, message broker, distributed caching, full-text search engine
% và identity provider vào một hệ thống vận hành thống nhất cũng là một thách thức kỹ thuật
% đáng kể.

% Xuất phát từ thực tiễn trên, đề tài đặt ra câu hỏi nghiên cứu: \textit{Làm thế nào để thiết
% kế và triển khai một hệ thống thương mại điện tử hoàn chỉnh theo kiến trúc microservices, có
% khả năng mở rộng độc lập theo từng miền nghiệp vụ, chịu lỗi cục bộ, đảm bảo tính nhất quán
% dữ liệu trong môi trường phân tán, và đặc biệt hỗ trợ tính năng flash sale chống over-sell
% trong điều kiện tải cao?}
Trong thập kỷ qua, thương mại điện tử (TMĐT) tại Việt Nam đã chứng kiến tốc độ tăng trưởng
vượt bậc. Theo báo cáo của Hiệp hội Thương mại điện tử Việt Nam (VECOM) năm 2024, doanh thu
TMĐT B2C ước đạt 25 tỷ USD, với mức tăng trưởng bình quân trên 25\%/năm trong giai đoạn
2020--2024. Các sàn giao dịch lớn như Shopee, Lazada, Tiki hiện phục vụ hàng chục triệu
người dùng đồng thời, đặc biệt trong các sự kiện Flash Sale và Black Friday --- thời điểm
lưu lượng giao dịch có thể tăng đột biến gấp 10--50 lần so với ngày thường. Quy mô và tính
biến động cao này đặt ra những yêu cầu khắt khe về mặt kiến trúc đối với các hệ thống TMĐT.

Thực tế vận hành cho thấy các hệ thống TMĐT xây dựng theo kiến trúc Monolith truyền thống
bộc lộ nhiều hạn chế nghiêm trọng khi đối mặt với bài toán quy mô lớn. Thứ nhất, toàn bộ hệ
thống buộc phải được mở rộng đồng thời ngay cả khi chỉ một module đang chịu tải cao, dẫn đến
lãng phí tài nguyên và gia tăng chi phí vận hành. Thứ hai, mọi thay đổi dù nhỏ đều yêu cầu
triển khai lại toàn bộ ứng dụng, làm chậm tốc độ phát hành tính năng mới và tăng rủi ro mỗi
khi cập nhật. Thứ ba, một lỗi cục bộ có thể kéo sập toàn bộ hệ thống do thiếu cơ chế cách ly
sự cố (fault isolation). Đối với những nghiệp vụ đặc thù như flash sale --- nơi hàng chục
nghìn người dùng đồng thời cạnh tranh một số lượng hàng giới hạn chỉ trong vài phút ---
kiến trúc Monolith gần như không thể đáp ứng yêu cầu về khả năng mở rộng chọn lọc và xử lý
đồng thời cao.

Kiến trúc microservices ra đời như một hướng tiếp cận nhằm khắc phục những hạn chế nêu trên,
bằng cách phân rã hệ thống thành các dịch vụ độc lập, có thể mở rộng và triển khai riêng
biệt. Tuy nhiên, lợi ích này không đến một cách miễn phí: việc chuyển từ một khối ứng dụng
duy nhất sang một hệ phân tán làm phát sinh hàng loạt thách thức mới về tính nhất quán dữ
liệu, độ tin cậy của truyền thông điệp và khả năng chịu lỗi giữa các dịch vụ. Giải quyết
trọn vẹn những thách thức này đòi hỏi vận dụng nhiều mẫu thiết kế phân tán phức tạp, đồng
thời tích hợp đồng bộ một loạt thành phần hạ tầng như service discovery, API gateway,
message broker, distributed caching, search engine và identity provider vào một hệ thống
vận hành thống nhất. Đây là những vấn đề được đề cập rộng rãi trong tài liệu học thuật, song
ít dự án thực tế tại Việt Nam triển khai một cách đầy đủ và có hệ thống.

Xuất phát từ thực tiễn trên, đề tài đặt ra câu hỏi nghiên cứu: \textit{Làm thế nào để thiết
kế và triển khai một hệ thống thương mại điện tử hoàn chỉnh theo kiến trúc microservices, có
khả năng mở rộng độc lập theo từng miền nghiệp vụ, chịu lỗi cục bộ, đảm bảo tính nhất quán
dữ liệu trong môi trường phân tán, và đặc biệt hỗ trợ tính năng flash sale chống over-sell
trong điều kiện tải cao?}

% \section{Mục tiêu đề tài}
% \label{section:1.2}

% \subsection*{Mục tiêu tổng quát}

% Thiết kế, triển khai và đánh giá một hệ thống thương mại điện tử đầy đủ chức năng theo kiến
% trúc microservices, sử dụng Spring Boot 3.5, Spring Cloud 2025 và Java 21, qua đó minh chứng
% tính khả thi và hiệu quả của kiến trúc này trong việc giải quyết các bài toán đặc thù của
% TMĐT: khả năng mở rộng chọn lọc, xử lý đồng thời cao, tính nhất quán dữ liệu phân tán và
% khả năng quan sát hệ thống từ đầu đến cuối.

% \subsection*{Mục tiêu cụ thể}

% Đề tài đặt ra sáu mục tiêu cụ thể, mỗi mục tiêu đều có thể đo lường và kiểm chứng thông qua
% mã nguồn, cấu hình hệ thống và kết quả kiểm thử:

% \begin{enumerate}
%     \item \textbf{Phân rã hệ thống theo Domain-Driven Design:} Phân tích và thiết kế 13
%     bounded context độc lập, mỗi context tương ứng với một microservice sở hữu cơ sở dữ
%     liệu, Dockerfile và vòng đời triển khai riêng biệt. Các service bao phủ đầy đủ các miền
%     nghiệp vụ cốt lõi của TMĐT: định danh, người dùng, sản phẩm, tồn kho, giỏ hàng, đơn
%     hàng, thanh toán, khuyến mãi, thông báo, đánh giá, tìm kiếm, nội dung và flash sale.

%     \item \textbf{Áp dụng các mẫu thiết kế phân tán cốt lõi:} Triển khai bốn mẫu thiết kế
%     quan trọng trong môi trường phân tán --- Saga Orchestration (với order-service làm bộ
%     điều phối giao dịch phân tán trong luồng đặt hàng); Transactional Outbox (đảm bảo phát
%     hành sự kiện at-least-once ngay cả khi service gặp sự cố); Idempotent Consumer (đảm
%     bảo hiệu ứng xử lý exactly-once ở phía nhận); và CQRS-lite (tách biệt mô hình đọc với
%     Elasticsearch phục vụ tìm kiếm toàn văn).

%     \item \textbf{Tích hợp bảo mật tập trung với Keycloak:} Sử dụng Keycloak 26 làm nhà cung
%     cấp định danh, cấu hình OAuth 2.0/OpenID Connect với JWT ký thuật toán RS256 (khóa bất
%     đối xứng). API Gateway đóng vai trò Resource Server duy nhất theo mẫu Trusted Subsystem
%     --- xác thực tập trung, chuyển tiếp thông tin định danh qua header tin cậy cho các
%     service backend, loại bỏ việc mỗi service phải gọi lại Keycloak để xác minh token.

%     \item \textbf{Tích hợp cổng thanh toán VNPAY:} Kết nối với môi trường sandbox của VNPAY,
%     sinh URL thanh toán với chữ ký xác thực HMAC-SHA512, xử lý Return URL và IPN callback,
%     đồng thời triển khai cơ chế tự động hủy thanh toán khi quá thời hạn 30 phút.

%     \item \textbf{Xây dựng cơ chế Flash Sale chống over-sell:} Thiết kế giải pháp chống tranh
%     chấp tài nguyên bằng Redis Lua atomic script, buyer set chống mua trùng, cơ chế bù trừ
%     khi tạo đơn thất bại và reconciliation scheduler đối chiếu lại số đơn thực tế. Mục tiêu:
%     trong điều kiện nhiều người dùng đồng thời, số đơn hàng thành công không vượt quá số
%     lượng slot của chiến dịch.

%     \item \textbf{Xây dựng hệ thống quan sát hoàn chỉnh:} Triển khai Prometheus và Grafana
%     cho thu thập và trực quan hóa số liệu (JVM heap, request rate, latency, error rate);
%     Zipkin cho truy vết phân tán theo chuẩn W3C Trace Context; structured logging với khả
%     năng liên kết traceId/spanId, cho phép giám sát và gỡ lỗi hệ thống từ đầu đến cuối.
% \end{enumerate}

\section{Mục tiêu đề tài}
\label{section:1.2}

\subsection*{Mục tiêu tổng quát}

Đề tài hướng tới thiết kế, triển khai và đánh giá một hệ thống thương mại điện tử đầy đủ chức
năng theo kiến trúc microservices, sử dụng ngôn ngữ Java 21, nền tảng Spring Boot 3.5 và hệ
sinh thái Spring Cloud 2025. Thông qua hệ thống này, đề tài minh chứng tính khả thi và hiệu
quả của kiến trúc microservices trong việc giải quyết các bài toán đặc thù của TMĐT đã nêu ở
phần Đặt vấn đề: khả năng mở rộng chọn lọc, xử lý đồng thời cao, đảm bảo tính nhất quán dữ
liệu trong môi trường phân tán và khả năng quan sát hệ thống từ đầu đến cuối.

\subsection*{Mục tiêu cụ thể}

Trên cơ sở mục tiêu tổng quát, đề tài xác định sáu mục tiêu cụ thể; mỗi mục tiêu đều có thể
đo lường và kiểm chứng thông qua mã nguồn, cấu hình hệ thống và kết quả kiểm thử:

\begin{enumerate}
    \item \textbf{Phân rã hệ thống theo Domain-Driven Design:} Phân tích và thiết kế các
    bounded context độc lập, mỗi context tương ứng với một microservice sở hữu cơ sở dữ liệu
    và vòng đời triển khai riêng biệt, bao phủ đầy đủ các miền nghiệp vụ cốt lõi của TMĐT.

    \item \textbf{Áp dụng các mẫu thiết kế phân tán cốt lõi:} Vận dụng bốn mẫu thiết kế quan
    trọng để giải quyết các thách thức của hệ phân tán --- Saga Orchestration cho giao dịch
    phân tán, Transactional Outbox cho truyền thông điệp tin cậy, Idempotent Consumer cho xử
    lý lặp ở phía nhận, và CQRS-lite tách biệt mô hình đọc phục vụ tìm kiếm toàn văn.

    \item \textbf{Tích hợp bảo mật tập trung với Keycloak:} Sử dụng Keycloak làm nhà cung cấp
    định danh theo chuẩn OAuth 2.0/OpenID Connect, với API Gateway đóng vai trò điểm xác thực
    tập trung theo mẫu Trusted Subsystem.

    \item \textbf{Tích hợp cổng thanh toán VNPAY:} Kết nối và xử lý hoàn chỉnh luồng thanh
    toán trực tuyến qua VNPAY, bao gồm cơ chế xác thực giao dịch và tự động hủy khi quá hạn.

    \item \textbf{Xây dựng cơ chế Flash Sale chống over-sell:} Thiết kế giải pháp đảm bảo
    trong điều kiện nhiều người dùng truy cập đồng thời, số đơn hàng thành công không vượt
    quá số lượng giới hạn của chiến dịch.

    \item \textbf{Xây dựng hệ thống quan sát hoàn chỉnh:} Triển khai khả năng quan sát trên
    cả ba trụ cột metrics, distributed tracing và structured logging, cho phép giám sát và
    gỡ lỗi hệ thống từ đầu đến cuối.
\end{enumerate}
% \section{Phạm vi đề tài}
% \label{section:1.3}

% \subsection*{Phạm vi thực hiện}

% Đề tài tập trung xây dựng phần backend của một hệ thống thương mại điện tử hoàn chỉnh theo
% kiến trúc microservices, kèm theo giao diện minh họa và bộ kịch bản kiểm thử thực nghiệm.
% Về quy mô, hệ thống gồm 13 business microservice tương ứng với 13 miền nghiệp vụ độc lập, ba
% Spring Cloud infrastructure service đảm nhiệm vai trò định tuyến --- phát hiện dịch vụ ---
% quản lý cấu hình, cùng các thành phần hạ tầng phục vụ lưu trữ, giao tiếp bất đồng bộ, tìm
% kiếm, xác thực và quan sát hệ thống. Về giao diện, đề tài triển khai một storefront cho phía
% khách hàng và một Admin Panel cho phía quản trị, đủ để minh họa các luồng nghiệp vụ chính.
% Toàn bộ hệ thống được đóng gói và vận hành bằng Docker Compose, đồng thời được triển khai ở
% mức production-like trên hạ tầng AWS. Bảng~\ref{table:in-scope} trình bày tổng quan các thành
% phần cùng vai trò của chúng trong hệ thống.

% \begin{table}[H]
% \centering
% \caption{Các thành phần trong phạm vi đề tài}
% \label{table:in-scope}
% \resizebox{\textwidth}{!}{%
% \begin{tabular}{|l|l|l|}
% \hline
% \textbf{Nhóm} & \textbf{Thành phần} & \textbf{Vai trò trong hệ thống} \\ \hline
% \multirow{3}{2.5cm}{Spring Cloud} & API Gateway & Điểm vào duy nhất, định tuyến \& xác thực \\ \cline{2-3}
%                                   & Discovery Server & Đăng ký \& phát hiện dịch vụ \\ \cline{2-3}
%                                   & Config Server & Quản lý cấu hình tập trung \\ \hline
% \multirow{13}{2.5cm}{Business Services} & identity-service, user-service & Định danh \& quản lý người dùng \\ \cline{2-3}
%                                         & product-service & Quản lý sản phẩm, danh mục \\ \cline{2-3}
%                                         & inventory-service & Quản lý tồn kho \\ \cline{2-3}
%                                         & cart-service & Giỏ hàng \\ \cline{2-3}
%                                         & order-service, payment-service & Đặt hàng \& thanh toán \\ \cline{2-3}
%                                         & voucher-service, content-service & Khuyến mãi \& nội dung \\ \cline{2-3}
%                                         & notification-service & Thông báo, gửi email \\ \cline{2-3}
%                                         & review-service & Đánh giá sản phẩm \\ \cline{2-3}
%                                         & search-service & Tìm kiếm toàn văn \\ \cline{2-3}
%                                         & flash-sale-service & Xử lý flash sale chống over-sell \\ \hline
% \multirow{5}{2.5cm}{Infrastructure} & PostgreSQL, Redis & Cơ sở dữ liệu \& cache \\ \cline{2-3}
%                                      & Kafka & Message broker bất đồng bộ \\ \cline{2-3}
%                                      & Elasticsearch & Search engine \\ \cline{2-3}
%                                      & Keycloak, Mailpit & Định danh tập trung \& email \\ \cline{2-3}
%                                      & Prometheus, Grafana, Zipkin & Khả năng quan sát hệ thống \\ \hline
% \multirow{2}{2.5cm}{Frontend} & Storefront & Giao diện mua sắm phía khách hàng \\ \cline{2-3}
%                               & Admin Panel & Trang quản trị (ROLE\_ADMIN) \\ \hline
% \multirow{4}{2.5cm}{Kiểm thử \& Triển khai} & Unit/Integration Test & Kiểm thử đơn vị \& tích hợp \\ \cline{2-3}
%                                             & Smoke/E2E Test & Kiểm thử luồng đầu cuối \\ \cline{2-3}
%                                             & Performance Test & Kiểm thử hiệu năng, tải cao \\ \cline{2-3}
%                                             & Triển khai & Docker Compose, AWS EC2 single-host \\ \hline
% \end{tabular}
% }
% \end{table}

% \subsection*{Ngoài phạm vi}

% Xuất phát từ giới hạn về thời gian, tài nguyên và định hướng nghiên cứu của một đồ án tốt
% nghiệp, các hạng mục sau được xác định nằm ngoài phạm vi đề tài:

% \begin{itemize}
%     \item \textbf{Frontend cấp sản phẩm thương mại:} Đề tài đã triển khai storefront và Admin
%     Panel phục vụ minh họa các luồng nghiệp vụ chính. Tuy nhiên, các năng lực ở mức production
%     như tối ưu trải nghiệm người dùng đa thiết bị, upload file qua object storage, workflow
%     kiểm duyệt nâng cao và ứng dụng mobile riêng chưa nằm trong phạm vi.

%     \item \textbf{Kubernetes và Service Mesh:} Việc triển khai trên Kubernetes kết hợp Istio
%     hoặc Linkerd được xem là hướng mở rộng trong tương lai. Phạm vi hiện tại dừng ở Docker
%     Compose single-host, phù hợp với quy mô của một đồ án tốt nghiệp.

%     \item \textbf{CI/CD pipeline cấp sản xuất đầy đủ:} Đề tài đã thiết lập quy trình build và
%     triển khai tự động ở mức production-like. Tuy nhiên, các năng lực như canary/blue-green
%     deployment, quét bảo mật bắt buộc, quản lý secret chuyên dụng và triển khai đa môi trường
%     chưa được đề cập.

%     \item \textbf{Thanh toán thực tế:} Đề tài chỉ sử dụng môi trường sandbox của cổng thanh
%     toán để minh họa luồng giao dịch, chưa đăng ký tài khoản merchant phục vụ thanh toán thật.

%     \item \textbf{Logistics và hoàn tiền:} Các nghiệp vụ vận chuyển, giao hàng và hoàn tiền
%     phức tạp nằm ngoài phạm vi backend e-commerce cốt lõi mà đề tài hướng tới.

%     \item \textbf{Triển khai đa vùng:} Hệ thống được triển khai trên một máy ảo duy nhất tại
%     một khu vực địa lý, chưa hỗ trợ cơ chế chuyển đổi dự phòng đa vùng.
% \end{itemize}

\section{Phạm vi đề tài}
\label{section:1.3}

\subsection*{Phạm vi thực hiện}

Đề tài tập trung xây dựng phần backend của một hệ thống thương mại điện tử hoàn chỉnh theo
kiến trúc microservices, kèm theo giao diện minh họa và bộ kịch bản kiểm thử thực nghiệm.
Về quy mô, hệ thống gồm 13 business microservice tương ứng với 13 miền nghiệp vụ độc lập, ba
Spring Cloud infrastructure service đảm nhiệm vai trò định tuyến, phát hiện dịch vụ và quản
lý cấu hình tập trung, cùng các thành phần hạ tầng phục vụ lưu trữ, giao tiếp bất đồng bộ,
tìm kiếm, xác thực và quan sát hệ thống. Về giao diện, đề tài triển khai một storefront cho
phía khách hàng và một Admin Panel cho phía quản trị, đủ để minh họa các luồng nghiệp vụ
chính. Toàn bộ hệ thống được đóng gói và vận hành bằng Docker Compose, đồng thời được triển
khai ở mức production-like trên hạ tầng AWS. Bảng~\ref{table:in-scope} trình bày tổng quan
các thành phần cùng vai trò của chúng trong hệ thống; chi tiết về lựa chọn và lý do sử dụng
từng công nghệ sẽ được phân tích trong Chương~3.

\begin{table}[H]
\centering
\caption{Các thành phần trong phạm vi đề tài}
\label{table:in-scope}
\resizebox{\textwidth}{!}{%
\begin{tabular}{|l|l|l|}
\hline
\textbf{Nhóm} & \textbf{Thành phần} & \textbf{Vai trò trong hệ thống} \\ \hline
\multirow{3}{2.5cm}{Spring Cloud} & API Gateway & Điểm vào duy nhất, định tuyến \& xác thực \\ \cline{2-3}
                                  & Discovery Server & Đăng ký \& phát hiện dịch vụ \\ \cline{2-3}
                                  & Config Server & Quản lý cấu hình tập trung \\ \hline
\multirow{13}{2.5cm}{Business Services} & identity-service, user-service & Định danh \& quản lý người dùng \\ \cline{2-3}
                                        & product-service & Quản lý sản phẩm, danh mục \\ \cline{2-3}
                                        & inventory-service & Quản lý tồn kho \\ \cline{2-3}
                                        & cart-service & Giỏ hàng \\ \cline{2-3}
                                        & order-service, payment-service & Đặt hàng \& thanh toán \\ \cline{2-3}
                                        & voucher-service, content-service & Khuyến mãi \& nội dung \\ \cline{2-3}
                                        & notification-service & Thông báo, gửi email \\ \cline{2-3}
                                        & review-service & Đánh giá sản phẩm \\ \cline{2-3}
                                        & search-service & Tìm kiếm toàn văn \\ \cline{2-3}
                                        & flash-sale-service & Xử lý flash sale chống over-sell \\ \hline
\multirow{5}{2.5cm}{Infrastructure} & PostgreSQL, Redis & Cơ sở dữ liệu \& cache \\ \cline{2-3}
                                     & Kafka & Message broker bất đồng bộ \\ \cline{2-3}
                                     & Elasticsearch & Search engine \\ \cline{2-3}
                                     & Keycloak, Mailpit & Định danh tập trung \& email \\ \cline{2-3}
                                     & Prometheus, Grafana, Zipkin & Khả năng quan sát hệ thống \\ \hline
\multirow{2}{2.5cm}{Frontend} & Storefront & Giao diện mua sắm phía khách hàng \\ \cline{2-3}
                              & Admin Panel & Trang quản trị (ROLE\_ADMIN) \\ \hline
\multirow{4}{2.5cm}{Kiểm thử \& Triển khai} & Unit/Integration Test & Kiểm thử đơn vị \& tích hợp \\ \cline{2-3}
                                            & Smoke/E2E Test & Kiểm thử luồng đầu cuối \\ \cline{2-3}
                                            & Performance Test & Kiểm thử hiệu năng, tải cao \\ \cline{2-3}
                                            & Triển khai & Docker Compose, AWS EC2 single-host \\ \hline
\end{tabular}
}
\end{table}

\subsection*{Ngoài phạm vi}

Xuất phát từ giới hạn về thời gian, tài nguyên và định hướng nghiên cứu của một đồ án tốt
nghiệp, các hạng mục sau được xác định nằm ngoài phạm vi đề tài:

\begin{itemize}
    \item \textbf{Frontend cấp sản phẩm thương mại:} Đề tài đã triển khai storefront và Admin
    Panel phục vụ minh họa các luồng nghiệp vụ chính. Tuy nhiên, các năng lực ở mức production
    như tối ưu trải nghiệm người dùng đa thiết bị, upload file qua object storage, workflow
    kiểm duyệt nâng cao và ứng dụng mobile riêng chưa nằm trong phạm vi.

    \item \textbf{Kubernetes và Service Mesh:} Việc triển khai trên Kubernetes kết hợp Istio
    hoặc Linkerd được xem là hướng mở rộng trong tương lai. Phạm vi hiện tại dừng ở Docker
    Compose single-host, phù hợp với quy mô của một đồ án tốt nghiệp.

    \item \textbf{CI/CD pipeline cấp sản xuất đầy đủ:} Đề tài đã thiết lập quy trình build và
    triển khai tự động ở mức production-like. Tuy nhiên, các năng lực như canary/blue-green
    deployment, quét bảo mật bắt buộc, quản lý secret chuyên dụng và triển khai đa môi trường
    chưa được đề cập.

    \item \textbf{Thanh toán thực tế:} Đề tài chỉ sử dụng môi trường sandbox của cổng thanh
    toán VNPAY để minh họa luồng giao dịch, chưa đăng ký tài khoản merchant phục vụ thanh
    toán thật.

    \item \textbf{Logistics và hoàn tiền:} Các nghiệp vụ vận chuyển, giao hàng và hoàn tiền
    phức tạp nằm ngoài phạm vi backend e-commerce cốt lõi mà đề tài hướng tới.

    \item \textbf{Triển khai đa vùng:} Hệ thống được triển khai trên một máy ảo duy nhất tại
    một khu vực địa lý, chưa hỗ trợ cơ chế chuyển đổi dự phòng đa vùng.
\end{itemize}
% \section{Khảo sát hệ thống tương tự}
% \label{section:1.4}

% Để có cơ sở thực tiễn cho các quyết định thiết kế, đề tài tiến hành khảo sát kiến trúc của
% một số nền tảng TMĐT lớn đang vận hành thành công trên thị trường. Mặc dù các nền tảng này
% vận hành theo mô hình sàn giao dịch (multi-vendor marketplace), các quyết định kỹ thuật của
% họ --- đặc biệt về kiến trúc microservices, chiến lược dữ liệu, message queue và cơ chế xử
% lý đồng thời cao --- là những bài học giá trị cho bất kỳ hệ thống TMĐT nào, bao gồm cả mô
% hình B2C mà đề tài hướng tới.

% \subsection*{Kiến trúc các nền tảng thương mại điện tử lớn}

% Bảng~\ref{table:survey-platforms} tổng hợp kết quả khảo sát kiến trúc của ba nền tảng TMĐT
% hàng đầu tại Việt Nam và Đông Nam Á, tập trung vào các khía cạnh: mô hình kiến trúc, ngôn
% ngữ lập trình chính, hệ quản trị cơ sở dữ liệu, message queue, giải pháp tìm kiếm và chiến
% lược xử lý flash sale.

% \begin{table}[H]
% \centering
% \caption{Kiến trúc các nền tảng thương mại điện tử lớn}
% \label{table:survey-platforms}
% \begin{tabular}{|l|p{3.2cm}|p{3.2cm}|p{3.2cm}|}
% \hline
% \textbf{Tiêu chí} & \textbf{Shopee} & \textbf{Tiki} & \textbf{Lazada} \\ \hline
% Kiến trúc & Microservices (hàng trăm service) & Microservices kết hợp Monolith & Microservices (hệ sinh thái Alibaba) \\ \hline
% Ngôn ngữ chính & Go, Python, Java & Java (Spring), Go & Java \\ \hline
% Cơ sở dữ liệu & MySQL sharding, Redis, TiDB & PostgreSQL, MySQL, Redis, ES & MySQL, Redis, HBase \\ \hline
% Message Queue & Kafka, RabbitMQ & Kafka & RocketMQ, Kafka \\ \hline
% Tìm kiếm & Elasticsearch & Elasticsearch & Elasticsearch \\ \hline
% Flash Sale & Redis + Lua + DB constraint & Redis + Queue + DB lock & Alibaba Sentinel \\ \hline
% \end{tabular}
% \end{table}

% Từ kết quả khảo sát, có thể rút ra năm điểm chung quan trọng trong kiến trúc kỹ thuật của
% các nền tảng TMĐT lớn, qua đó củng cố cho các quyết định công nghệ của đề tài: (1)~kiến
% trúc microservices được sử dụng để phân rã nghiệp vụ phức tạp thành các dịch vụ độc lập;
% (2)~chiến lược polyglot persistence --- kết hợp nhiều loại cơ sở dữ liệu cho các mục đích
% khác nhau --- là cách tiếp cận phổ biến; (3)~message queue đóng vai trò trung tâm trong
% xử lý bất đồng bộ và giao tiếp giữa các dịch vụ; (4)~Elasticsearch là lựa chọn thống nhất
% cho bài toán tìm kiếm toàn văn; và (5)~Redis là thành phần then chốt trong kiến trúc xử lý
% flash sale, thường được kết hợp với các cơ chế bảo vệ ở tầng cơ sở dữ liệu.

% \subsection*{Tài liệu và dự án tham khảo}

% Bên cạnh việc khảo sát xu hướng kiến trúc, đề tài tham khảo các tài liệu và dự án sau
% trong quá trình thiết kế và triển khai:

% \begin{itemize}
%     \item \textbf{Microservices Patterns (Chris Richardson):} Đây là tài liệu tham khảo
%     nền tảng về các mẫu thiết kế trong kiến trúc microservices, bao gồm Saga, Transactional
%     Outbox, CQRS, API Gateway và Database per Service. Trang web microservices.io cung cấp
%     bộ pattern catalog đầy đủ kèm ví dụ minh họa, là nguồn tham khảo chính cho các quyết
%     định thiết kế phân tán trong đề tài.

%     \item \textbf{Building Microservices (Sam Newman, O'Reilly):} Cuốn sách kinh điển về
%     thiết kế và triển khai microservices, cung cấp nền tảng lý thuyết vững chắc về
%     bounded context, database per service, giao tiếp đồng bộ/bất đồng bộ, resilience
%     pattern và chiến lược triển khai. Ấn bản thứ hai (2021) cập nhật các thực tiễn mới
%     về container hóa và observability.

%     \item \textbf{Domain-Driven Design (Eric Evans):} Cuốn sách đặt nền móng cho phương
%     pháp DDD, cung cấp các khái niệm cốt lõi như Bounded Context, Ubiquitous Language,
%     Aggregate và Context Map --- những công cụ được sử dụng trực tiếp để phân rã hệ thống
%     thành 13 microservice trong đề tài.

%     \item \textbf{Spring Boot \& Spring Cloud Official Documentation:} Tài liệu chính thức
%     của Spring cung cấp hướng dẫn chi tiết về cách cấu hình Spring Cloud Gateway, Eureka
%     Service Discovery, Config Server, OpenFeign, Spring Kafka và Spring Security OAuth2
%     Resource Server. Đây là nguồn tham khảo kỹ thuật chính trong quá trình cài đặt.

%     \item \textbf{Các dự án e-commerce mã nguồn mở trên GitHub:} Nhiều dự án e-commerce
%     mã nguồn mở như Saleor, Medusa và Spree cung cấp ý tưởng tham khảo về cấu trúc thư
%     mục, thiết kế REST API và tổ chức dữ liệu. Tuy nhiên, hầu hết các dự án này chưa
%     triển khai đầy đủ các mẫu thiết kế phân tán phức tạp như Saga Orchestration,
%     Transactional Outbox với Kafka, Idempotent Consumer hay cơ chế Redis Lua atomic
%     counter cho flash sale --- những điểm mà đề tài này hướng tới giải quyết một cách
%     trọn vẹn.
% \end{itemize}

\section{Khảo sát hệ thống tương tự}
\label{section:1.4}

Để có cơ sở thực tiễn cho các quyết định thiết kế, đề tài tiến hành khảo sát kiến trúc của
một số nền tảng TMĐT lớn đang vận hành thành công trên thị trường. Mặc dù các nền tảng này
hoạt động theo mô hình sàn giao dịch (multi-vendor marketplace), những quyết định kỹ thuật
của họ --- đặc biệt về kiến trúc microservices, chiến lược dữ liệu, message queue và cơ chế
xử lý đồng thời cao --- vẫn là các bài học giá trị cho bất kỳ hệ thống TMĐT nào, bao gồm cả
mô hình B2C mà đề tài hướng tới.

\subsection*{Kiến trúc các nền tảng thương mại điện tử lớn}

Bảng~\ref{table:survey-platforms} tổng hợp kết quả khảo sát kiến trúc của ba nền tảng TMĐT
hàng đầu tại Việt Nam và Đông Nam Á, tập trung vào các khía cạnh: mô hình kiến trúc, ngôn
ngữ lập trình chính, hệ quản trị cơ sở dữ liệu, message queue, giải pháp tìm kiếm và chiến
lược xử lý flash sale.

\begin{table}[H]
\centering
\caption{Kiến trúc các nền tảng thương mại điện tử lớn}
\label{table:survey-platforms}
\begin{tabular}{|l|p{3.2cm}|p{3.2cm}|p{3.2cm}|}
\hline
\textbf{Tiêu chí} & \textbf{Shopee} & \textbf{Tiki} & \textbf{Lazada} \\ \hline
Kiến trúc & Microservices (hàng trăm service) & Microservices kết hợp Monolith & Microservices (hệ sinh thái Alibaba) \\ \hline
Ngôn ngữ chính & Go, Python, Java & Java (Spring), Go & Java \\ \hline
Cơ sở dữ liệu & MySQL sharding, Redis, TiDB & PostgreSQL, MySQL, Redis, ES & MySQL, Redis, HBase \\ \hline
Message Queue & Kafka, RabbitMQ & Kafka & RocketMQ, Kafka \\ \hline
Tìm kiếm & Elasticsearch & Elasticsearch & Elasticsearch \\ \hline
Flash Sale & Redis + Lua + DB constraint & Redis + Queue + DB lock & Alibaba Sentinel \\ \hline
\end{tabular}
\end{table}

Từ kết quả khảo sát, có thể rút ra năm điểm chung quan trọng trong kiến trúc kỹ thuật của
các nền tảng TMĐT lớn: (1)~kiến trúc microservices được sử dụng để phân rã nghiệp vụ phức
tạp thành các dịch vụ độc lập; (2)~chiến lược polyglot persistence --- kết hợp nhiều loại
cơ sở dữ liệu cho các mục đích khác nhau --- là cách tiếp cận phổ biến; (3)~message queue
đóng vai trò trung tâm trong xử lý bất đồng bộ và giao tiếp giữa các dịch vụ; (4)~Elasticsearch
là lựa chọn thống nhất cho bài toán tìm kiếm toàn văn; và (5)~Redis là thành phần then chốt
trong kiến trúc xử lý flash sale, thường được kết hợp với các cơ chế bảo vệ ở tầng cơ sở dữ
liệu. Năm bài học này là cơ sở thực tiễn quan trọng, định hướng trực tiếp cho các quyết định
công nghệ được trình bày ở phần Định hướng giải pháp tiếp theo.

% \subsection*{Tài liệu và dự án tham khảo}

% Bên cạnh việc khảo sát xu hướng kiến trúc, đề tài tham khảo các tài liệu và dự án sau
% trong quá trình thiết kế và triển khai:

% \begin{itemize}
%     \item \textbf{Microservices Patterns (Chris Richardson):} Tài liệu tham khảo nền tảng về
%     các mẫu thiết kế trong kiến trúc microservices, bao gồm Saga, Transactional Outbox, CQRS,
%     API Gateway và Database per Service. Trang web microservices.io cung cấp bộ pattern catalog
%     đầy đủ kèm ví dụ minh họa, là nguồn tham khảo chính cho các quyết định thiết kế phân tán
%     trong đề tài.

%     \item \textbf{Building Microservices (Sam Newman, O'Reilly):} Cuốn sách kinh điển về
%     thiết kế và triển khai microservices, cung cấp nền tảng lý thuyết vững chắc về bounded
%     context, database per service, giao tiếp đồng bộ/bất đồng bộ, resilience pattern và
%     chiến lược triển khai. Ấn bản thứ hai (2021) cập nhật các thực tiễn mới về container hóa
%     và observability.

%     \item \textbf{Domain-Driven Design (Eric Evans):} Cuốn sách đặt nền móng cho phương pháp
%     DDD, cung cấp các khái niệm cốt lõi như Bounded Context, Ubiquitous Language, Aggregate
%     và Context Map --- những công cụ được vận dụng trực tiếp để phân rã hệ thống thành các
%     microservice trong đề tài.

%     \item \textbf{Spring Boot \& Spring Cloud Official Documentation:} Tài liệu chính thức
%     của Spring cung cấp hướng dẫn chi tiết về cách cấu hình Spring Cloud Gateway, Eureka
%     Service Discovery, Config Server, OpenFeign, Spring Kafka và Spring Security OAuth2
%     Resource Server. Đây là nguồn tham khảo kỹ thuật chính trong quá trình cài đặt.

%     \item \textbf{Các dự án e-commerce mã nguồn mở trên GitHub:} Nhiều dự án e-commerce mã
%     nguồn mở như Saleor, Medusa và Spree cung cấp ý tưởng tham khảo về cấu trúc thư mục,
%     thiết kế REST API và tổ chức dữ liệu. Tuy nhiên, hầu hết các dự án này chưa triển khai
%     đầy đủ các mẫu thiết kế phân tán phức tạp như Saga Orchestration, Transactional Outbox
%     với Kafka, Idempotent Consumer hay cơ chế Redis Lua atomic counter cho flash sale ---
%     những điểm mà đề tài này hướng tới giải quyết một cách trọn vẹn.
% \end{itemize}
% % \section{Định hướng giải pháp}
% % \label{section:1.5}

% % Dựa trên phân tích bài toán, khảo sát thực tiễn và nghiên cứu lý thuyết, đề tài xác định
% % sáu định hướng kỹ thuật then chốt, chi phối toàn bộ quá trình thiết kế và triển khai hệ
% % thống.

% % \begin{figure}[H]
% %     \centering
% %     \includegraphics[width=0.85\textwidth]{Hinhve/chuong1/6-tru-cot-ki-thuat.png}
% %     \caption{Sơ đồ tổng quan sáu trụ cột kỹ thuật của hệ thống}
% %     \label{fig:solution-overview}
% % \end{figure}

% % \textbf{Kiến trúc và phân rã dịch vụ.} Áp dụng phương pháp Domain-Driven Design với Bounded
% % Context làm đơn vị phân rã, hệ thống được chia thành 13 microservice tương ứng với 13 miền
% % nghiệp vụ độc lập. Mỗi service sở hữu cơ sở dữ liệu riêng và không được phép truy cập trực
% % tiếp vào dữ liệu của service khác --- mọi tương tác dữ liệu phải thông qua API hoặc sự kiện.
% % API Gateway đóng vai trò điểm vào duy nhất cho mọi yêu cầu từ phía client, kết hợp với
% % Eureka Service Discovery để định tuyến động và Config Server để quản lý cấu hình tập trung,
% % đảm bảo tính nhất quán về mặt vận hành trên toàn hệ thống.

% % \textbf{Giao tiếp và tính nhất quán phân tán.} Hệ thống kết hợp hai phương thức giao tiếp:
% % REST + OpenFeign cho các truy vấn đồng bộ cần phản hồi tức thời (xác minh giá sản phẩm,
% % kiểm tra tồn kho) và Apache Kafka KRaft cho các tương tác bất đồng bộ không yêu cầu phản hồi
% % ngay (cập nhật trạng thái đơn hàng, gửi thông báo). Saga Orchestration Pattern được áp dụng
% % để quản lý giao dịch phân tán trong luồng đặt hàng, với order-service đóng vai trò bộ điều
% % phối trung tâm, đảm bảo tính nhất quán cuối cùng thay vì two-phase commit. Toàn bộ 13 Kafka
% % topic được thiết kế rõ ràng về producer, consumer và cấu trúc payload.

% % \textbf{Truyền thông điệp tin cậy.} Transactional Outbox Pattern được triển khai tại
% % order-service và payment-service để giải quyết bài toán dual-write: sự kiện được ghi vào
% % bảng \texttt{outbox} hoặc \texttt{payment\_outbox} trong cùng một giao dịch cơ sở dữ liệu với
% % dữ liệu nghiệp vụ. OutboxPoller (chu kỳ quét 1 giây) định kỳ đọc các sự kiện chưa được phát
% % hành bằng truy vấn \texttt{FOR UPDATE SKIP LOCKED} và gửi lên Kafka, đảm bảo at-least-once
% % delivery ngay cả khi service gặp sự cố giữa chừng. Phía nhận, Idempotent Consumer sử dụng
% % bảng \texttt{processed\_events} để phát hiện và loại bỏ sự kiện trùng lặp, đạt được hiệu ứng
% % exactly-once processing ở tầng ứng dụng. Sáu scheduled job chạy nền đảm bảo hệ thống tự phục
% % hồi sau các tình huống bất thường: ReservationExpiryScheduler, PaymentTimeoutScheduler,
% % OutboxPoller cho order và payment, CampaignScheduler và ReconciliationScheduler.

% % \textbf{Flash sale chống over-sell.} Giải pháp được thiết kế theo mô hình phòng thủ đa lớp:
% % lớp thứ nhất sử dụng Redis Lua atomic script để thực hiện thao tác kiểm tra buyer set, kiểm
% % tra tồn kho và giảm slot trong bộ nhớ một cách nguyên tử; lớp thứ hai sử dụng logic tạo đơn
% % và cơ chế bù trừ để hoàn slot khi quá trình tạo đơn thất bại; lớp thứ ba sử dụng
% % ReconciliationScheduler để đối chiếu \texttt{soldCount} với số đơn flash-sale thực tế trong
% % order-service. Cơ chế distributed lock dựa trên Redis đảm bảo chỉ một instance của
% % flash-sale-service thực thi CampaignScheduler tại mỗi thời điểm, còn Redis stock key và
% % buyers key được đặt TTL theo thời điểm kết thúc chiến dịch.

% % \textbf{Bảo mật tập trung.} Toàn bộ hoạt động xác thực và phân quyền được xử lý tại API
% % Gateway thông qua Keycloak 26 theo mẫu Trusted Subsystem Pattern: Gateway là điểm duy nhất
% % xác minh JWT (RS256), sau đó chuyển tiếp thông tin định danh đã được xác thực qua ba header
% % tin cậy --- \texttt{X-User-Id}, \texttt{X-User-Roles} và \texttt{X-User-Email} --- cho các
% % service backend. Các service nội bộ hoàn toàn tin tưởng các header này mà không cần gọi lại
% % Keycloak, giúp giảm độ trễ và giảm tải cho Identity Provider. Phân quyền dựa trên mô hình
% % RBAC với hai vai trò chính: ROLE\_USER (khách hàng) và ROLE\_ADMIN (quản trị viên). Các
% % endpoint nhạy cảm được bảo vệ bổ sung bằng Redis token bucket rate limiter: 10 request/giây
% % cho đặt hàng và 3 request/giây cho mua flash sale.

% % \textbf{Khả năng quan sát.} Hệ thống quan sát được xây dựng trên cả ba trụ cột. Về
% % metrics: Prometheus định kỳ thu thập số liệu từ các endpoint \path|/actuator/prometheus|,
% % Grafana trực quan hóa thông qua các dashboard được provision tự động cho JVM, Spring Boot
% % service metrics và tổng quan Saga.
% % Về distributed tracing: Zipkin tiếp nhận và hiển thị vết (trace) từ tất cả service theo
% % chuẩn W3C Trace Context, mỗi request được gán một traceId duy nhất và mỗi đoạn xử lý trong
% % một service là một span, cho phép tái tạo toàn bộ hành trình của request qua hệ thống. Về
% % logging: structured logging với Logback pattern chứa traceId và spanId trong mỗi dòng log,
% % cho phép liên kết chéo giữa logs và traces khi cần điều tra sự cố.

\section{Định hướng giải pháp}
\label{section:1.5}

Dựa trên phân tích bài toán, khảo sát thực tiễn và nghiên cứu lý thuyết, đề tài xác định
sáu định hướng kỹ thuật then chốt, chi phối toàn bộ quá trình thiết kế và triển khai hệ
thống.

\begin{figure}[H]
    \centering
    \includegraphics[width=0.85\textwidth]{Hinhve/chuong1/6-tru-cot-ki-thuat.png}
    \caption{Sơ đồ tổng quan sáu trụ cột kỹ thuật của hệ thống}
    \label{fig:solution-overview}
\end{figure}

\textbf{Kiến trúc và phân rã dịch vụ.} Đề tài áp dụng phương pháp Domain-Driven Design với
Bounded Context làm đơn vị phân rã, chia hệ thống thành các microservice tương ứng với các
miền nghiệp vụ độc lập. Mỗi service sở hữu cơ sở dữ liệu riêng theo nguyên tắc Database per
Service và không truy cập trực tiếp dữ liệu của service khác --- mọi tương tác phải thông qua
API hoặc sự kiện. API Gateway đóng vai trò điểm vào duy nhất, phối hợp với Service Discovery
và Config Server để định tuyến động và quản lý cấu hình tập trung, đảm bảo tính nhất quán về
mặt vận hành trên toàn hệ thống.

\textbf{Giao tiếp và tính nhất quán phân tán.} Hệ thống kết hợp hai phương thức giao tiếp:
giao tiếp đồng bộ qua REST cho các truy vấn cần phản hồi tức thời (xác minh giá sản phẩm,
kiểm tra tồn kho) và giao tiếp bất đồng bộ qua Apache Kafka cho các tương tác không yêu cầu
phản hồi ngay (cập nhật trạng thái đơn hàng, gửi thông báo). Để quản lý giao dịch trải dài
trên nhiều service trong luồng đặt hàng, đề tài lựa chọn Saga Orchestration Pattern thay cho
two-phase commit, đảm bảo tính nhất quán cuối cùng (eventual consistency) mà không khóa tài
nguyên xuyên service.

\textbf{Truyền thông điệp tin cậy.} Để giải quyết bài toán dual-write --- nguy cơ mất đồng
bộ giữa việc ghi dữ liệu nghiệp vụ và phát hành sự kiện --- đề tài áp dụng Transactional
Outbox Pattern, ghi sự kiện vào bảng outbox trong cùng giao dịch với dữ liệu nghiệp vụ, qua
đó bảo đảm at-least-once delivery ngay cả khi service gặp sự cố. Ở phía nhận, Idempotent
Consumer Pattern loại bỏ các sự kiện trùng lặp, đạt hiệu ứng xử lý exactly-once ở tầng ứng
dụng. Bên cạnh đó, hệ thống sử dụng các tác vụ nền (scheduled job) để tự động phục hồi sau
những tình huống bất thường như sự kiện chưa được phát hành hay đặt chỗ tồn kho quá hạn.

\textbf{Flash sale chống over-sell.} Đây là bài toán đặc thù và thách thức nhất của đề tài,
được giải quyết theo mô hình phòng thủ đa lớp. Lớp thứ nhất tận dụng Redis kết hợp Lua script
để thực hiện kiểm tra điều kiện và trừ tồn kho một cách nguyên tử trong bộ nhớ, triệt tiêu
race condition ngay tại điểm nóng. Lớp thứ hai áp dụng cơ chế bù trừ để hoàn lại slot khi quá
trình tạo đơn thất bại. Lớp thứ ba dùng tác vụ đối chiếu định kỳ để phát hiện và vá các sai
lệch dữ liệu còn sót, bảo đảm số đơn thành công luôn khớp với số lượng giới hạn của chiến dịch.

\textbf{Bảo mật tập trung.} Toàn bộ hoạt động xác thực và phân quyền được tập trung tại API
Gateway theo mẫu Trusted Subsystem Pattern: Gateway là điểm duy nhất xác minh JWT do Keycloak
phát hành, sau đó chuyển tiếp thông tin định danh đã xác thực qua các header tin cậy cho các
service backend. Nhờ đó, các service nội bộ không phải lặp lại thao tác xác minh token, giúp
giảm độ trễ và giảm tải cho Identity Provider. Việc phân quyền dựa trên mô hình RBAC, kết hợp
với cơ chế rate limiting nhằm bảo vệ các endpoint nhạy cảm trước tình trạng quá tải.

\textbf{Khả năng quan sát.} Hệ thống quan sát được xây dựng trên cả ba trụ cột. Về metrics,
Prometheus thu thập số liệu vận hành và Grafana trực quan hóa qua các dashboard. Về
distributed tracing, Zipkin theo dõi hành trình của mỗi request xuyên suốt các service theo
chuẩn W3C Trace Context, gán một định danh truy vết (traceId) duy nhất cho từng request. Về
logging, cơ chế structured logging gắn traceId vào mỗi dòng log, cho phép liên kết chéo giữa
log và trace khi điều tra sự cố. Ba trụ cột này phối hợp tạo nên khả năng giám sát và gỡ lỗi
hệ thống từ đầu đến cuối.
% \section{Bố cục báo cáo}
% \label{section:1.6}

% Báo cáo được tổ chức thành năm chương chính, một chương kết luận và hai phụ lục. Nội dung
% các chương được sắp xếp theo logic từ tổng quan đến chi tiết, từ lý thuyết đến thực hành,
% từ thiết kế đến kiểm chứng:

% \textbf{Chương 1 --- Giới thiệu đề tài} trình bày bối cảnh, câu hỏi nghiên cứu, mục tiêu,
% phạm vi, khảo sát hệ thống tương tự và định hướng giải pháp.

% \textbf{Chương 2 --- Cơ sở lý thuyết} trình bày nền tảng lý thuyết về kiến trúc
% microservices, các mẫu thiết kế phân tán (Saga, Outbox, Idempotent Consumer, CQRS),
% cơ chế đảm bảo độ tin cậy, bảo mật phân tán và khả năng quan sát.

% \textbf{Chương 3 --- Giải pháp công nghệ} phân tích và lý giải việc lựa chọn từng công
% nghệ trong hệ thống, từ nền tảng phát triển (Java 21, Spring Boot, Spring Cloud), hạ
% tầng dữ liệu (PostgreSQL, Redis, Elasticsearch, Kafka), đến bảo mật (Keycloak), thanh
% toán (VNPAY), giao diện (Next.js), triển khai (Docker, AWS, CI/CD) và công cụ kiểm thử.

% \textbf{Chương 4 --- Phân tích và Thiết kế hệ thống} trình bày kết quả phân tích yêu cầu
% và thiết kế hệ thống, bao gồm kiến trúc tổng thể, thiết kế cơ sở dữ liệu, API, luồng sự
% kiện, bảo mật, máy trạng thái và các luồng nghiệp vụ chính.

% \textbf{Chương 5 --- Cài đặt, Triển khai, Kiểm thử và Đánh giá} mô tả quá trình hiện thực
% hóa thiết kế và triển khai production-like trên AWS EC2, đồng thời trình bày kết quả kiểm
% thử chức năng, hiệu năng và khả năng chịu lỗi của hệ thống.

% \textbf{Kết luận và Hướng phát triển} tổng kết kết quả đạt được, phân tích hạn chế và đề
% xuất các hướng phát triển trong tương lai.

% \textbf{Phụ lục A} cung cấp hướng dẫn cài đặt và vận hành hệ thống. \textbf{Phụ lục B}
% trình bày đặc tả chi tiết các use case chính.

\section{Bố cục báo cáo}
\label{section:1.6}

Báo cáo được tổ chức thành năm chương chính, một chương kết luận và hai phụ lục, sắp xếp
theo logic từ tổng quan đến chi tiết, từ lý thuyết đến thực hành và từ thiết kế đến kiểm chứng:

\begin{itemize}
    \item \textbf{Chương 1 --- Giới thiệu đề tài:} trình bày bối cảnh, câu hỏi nghiên cứu,
    mục tiêu, phạm vi, khảo sát hệ thống tương tự và định hướng giải pháp.

    \item \textbf{Chương 2 --- Cơ sở lý thuyết:} trình bày nền tảng lý thuyết về kiến trúc
    microservices, các mẫu thiết kế phân tán (Saga, Outbox, Idempotent Consumer, CQRS), bảo
    mật phân tán và khả năng quan sát.

    \item \textbf{Chương 3 --- Giải pháp công nghệ:} phân tích và lý giải việc lựa chọn từng
    công nghệ trong hệ thống, từ nền tảng phát triển, hạ tầng dữ liệu, bảo mật, thanh toán,
    giao diện đến triển khai và kiểm thử.

    \item \textbf{Chương 4 --- Phân tích và Thiết kế hệ thống:} trình bày kết quả phân tích
    yêu cầu và thiết kế hệ thống, bao gồm kiến trúc tổng thể, cơ sở dữ liệu, API, luồng sự
    kiện, bảo mật, máy trạng thái và các luồng nghiệp vụ chính.

    \item \textbf{Chương 5 --- Cài đặt, Triển khai và Kiểm thử :} mô tả quá trình
    hiện thực hóa thiết kế, triển khai production, đồng thời trình bày kết quả
    kiểm thử chức năng, hiệu năng và khả năng chịu lỗi.

    \item \textbf{Kết luận và Hướng phát triển:} tổng kết kết quả đạt được, phân tích hạn chế
    và đề xuất các hướng phát triển trong tương lai.
\end{itemize}

% Ngoài ra, \textbf{Phụ lục A} cung cấp hướng dẫn cài đặt và vận hành hệ thống, còn
% \textbf{Phụ lục B} trình bày đặc tả chi tiết các use case chính.
\end{document}
